\chapter{The cache controller}\label{ch:cache-controller}

\whatyoufind{What every cache in Octopus is made of, whether it is a private L1 or the shared
LLC: two interfaces, one processing queue, a protocol handler that reads a CSV state machine,
and a data handler that owns the array, the MSHRs and the write-back buffer. How a message
becomes a list of actions; what the controller admits and what it stalls; how the fields of a
\code{Message} are decoded; and the one ordering rule at the bus interface that a snoop
protocol's correctness rests on.}

\section{One class, every level}

A cache controller is the same object at every level of the hierarchy. \code{BaseController}
holds the machinery; \code{CacheController} adds the data array, the miss handling and the
write-back path; \code{CacheControllerExclusive} (the shipped L1), \code{CacheController\_End2End}
(the shipped LLC) and \code{CacheControllerDirectory} (the directory home) specialise only what
differs. What makes one instance an L1 and another the LLC is its position in the system and the
FSM table it is given, not its class.

\figmed{CacheController3.pdf}{Inside a cache controller. Messages from both interfaces are
serialised into one processing queue; the protocol handler turns each one into actions by
looking the (state, event) pair up in the FSM table that the FSM reader parsed at start-up;
the data handler executes those actions against the array, the MSHRs and the write-back
buffer, arbitrating the data port.}{fig:controller}{0.62}

\paragraph{Two interfaces.} Every controller has a \emph{lower} \code{CommunicationInterface},
facing the cores, and an \emph{upper} one, facing the interconnect towards memory. Demand
requests arrive from below; coherence requests, snooped transactions and data arrive from
above; the controller's own requests leave through the upper interface. Interfaces are the
only way in or out (\cref{ch:system-organization}): they are bounded queues, so a full
interface pushes back on whoever is trying to fill it.

\paragraph{The processing queue.} Messages from both interfaces are serialised into one
queue ordered \emph{first-ready, first-come-first-served} (\code{FRFCFS\_Buffer}). A message
is \emph{ready} when the protocol handler judges that its line is in a state that can accept
it; a message whose line is in a transient state waits without blocking the ones behind it
that can proceed. The queue is bounded by \key{processing\_queue\_size}, and the bound applies
to demand admission only: maintenance traffic the controller generates for itself, such as a
write-back draining the buffer, is admitted with \code{force=true} so it can never be
deadlocked by the demand it is making room for.

\paragraph{Per-line order.} Among messages to the \emph{same} line the queue keeps strict
arrival order for demand requests and invalidations together; data messages are exempt, so a
fill or a supply can always be processed and can never wait behind a stalled request. This
per-line FCFS gate is what prevents an eviction's invalidation from leapfrogging an older
\code{GetS}/\code{GetM} to the same line, and the Logger reports when a request was held by it
(\col{LLC Gate}, \cref{ch:monitoring}).

\section{From message to actions}

\begin{figure}[htbp]\centering
% One message through a cache controller: queue -> protocol handler -> FSMReader -> actions.
\begin{sequencediagram}
  \tikzstyle{inststyle}+=[rounded corners=2pt, fill=octmem, draw=octrule, font=\sffamily\small]
  \newthread[octctl]{q}{Processing queue}
  \newinst[1.6]{ph}{Protocol handler}
  \newinst[1.6]{fsm}{FSMReader}
  \newinst[1.8]{dh}{Data handler, interfaces}
  \begin{call}{q}{\code{processRequest(msg)}}{ph}{\code{ControllerAction[]}}
    \begin{callself}{ph}{\code{readEvent(msg)} $\rightarrow$ \code{EventId}}{}
    \end{callself}
    \begin{call}{ph}{\code{getTransition(state, event)}}{fsm}{\code{next\_state, actions[]}}
    \end{call}
    \begin{callself}{ph}{\code{handleAction(actions)}}{}
    \end{callself}
  \end{call}
  \begin{messcall}{q}{\code{SEND\_BUS\_MSG} / \code{UPDATE\_CACHE\_LINE} / \code{WRITE\_BACK} / \ldots}{dh}
  \end{messcall}
\end{sequencediagram}

\caption{One message through the controller. The handler decodes the message into an event,
looks up the transition for the line's current state, and hands the controller a list of
protocol-agnostic actions; the controller does not know which protocol it is running.}
\label{fig:controller-sequence}
\end{figure}

The protocol handler (\code{CoherenceProtocolHandler}, specialised per family:
\code{MSIProtocol}, \code{MESIProtocol}, \code{MOESIProtocol}, their \code{LLC\ldots} and
\code{\ldots Directory} variants) does three things for every message
(\cref{fig:controller-sequence}):

\begin{enumerate}
\item \textbf{Decode the event.} \code{readEvent} maps the message to one column of the FSM
  table from its \emph{source} and \emph{contents}: a \code{Load}/\code{Store} from below; an
  \code{Own\_GetS}/\code{Own\_GetM} when the controller sees its own request echoed back from the
  bus and \code{Other\_GetS}/\code{Other\_GetM} when it is another core's; \code{OwnData} when the
  message carries data for a line it is waiting on; \code{Invalidation} for an LLC
  back-invalidation; \code{Replacement} for the controller's own eviction.
\item \textbf{Look up the transition.} \code{FSMReader::getTransition(state, event)} returns the
  actions and next state the CSV table lists for that cell (\cref{ch:coherence-engine}).
\item \textbf{Translate to controller actions.} \code{handleAction} turns protocol actions such
  as \act{GetS}, \act{Data2Both} or \act{IssueInv} into \code{ControllerAction}s the controller
  can execute without knowing the protocol: \code{SEND\_BUS\_MSG}, \code{UPDATE\_CACHE\_LINE},
  \code{WRITE\_BACK}, \code{HIT\_Action}, \code{REMOVE\_PENDING}, \ldots
\end{enumerate}

A cell that reads \code{Fault/} means ``this (state, event) must be unreachable''; the handler
halts the run and reports it rather than mis-routing silently. That is what makes a protocol
auditable: an incomplete table fails loudly, on the first trace that exercises the hole.

\section{The data handler}

\code{CacheDataHandler\_COTS} owns everything that holds a line:

\begin{itemize}
\item \textbf{The array.} One array per controller, indexed by set, holding each line's tag,
  coherence bits and data together. Reading a line's \emph{bits} (the tag compare, the state
  the FSM needs) is untimed and claims nothing. Reading or writing its \emph{data} goes through a
  single port that is busy for \key{m\_data\_handler.m\_data\_access\_latency} cycles per access; demand
  reads, eviction reads and fill or write-back writes contend for it through a round-robin
  arbiter (\code{m\_data\_access\_arbiter}).
\item \textbf{The MSHRs.} A miss allocates one entry and waits there, not in the array, until
  its data returns. \key{num\_mshr} bounds the number of outstanding misses; \code{-1} removes
  the bound.
\item \textbf{The write-back buffer.} A victim moves out of the array into the buffer, bits and
  data, the moment its way is needed, and leaves for memory later as a separate transaction.
  \key{m\_data\_handler.pwb\_size} bounds it.
\item \textbf{Replacement.} LRU, FIFO or random, chosen per controller; under way partitioning
  (\cref{ch:tasks-and-demo}) a core's fills install into, and evict from, its own ways only.
\end{itemize}

A line that currently sits in the MSHR or the write-back buffer is read without waiting for the
port, because it is not in the array yet; it still restarts the port's busy timer, which is why
such accesses cut in ahead of an access already in progress (the dotted bypass in
\cref{fig:llc-architecture}). Fill writes are not on the requester's critical path: the
response is emitted from the fill message itself and the array write runs afterwards
(\cref{ch:llc-datapath,ch:request-end-to-end}).

\subsection{Two array models, and the line interlock}

The data port can be timed in two ways, chosen per controller
(\cref{tab:array-models}). In the \emph{occupancy} model, the default, a row runs when its
message is popped --- the line's state changes at once --- and only the access to the line's
bytes waits for the port, which then stays busy for the latency; the requester is not charged
the latency and one access is served per latency window. It is the model every standalone
preset uses, and it is correct for them because a trace carries no data and the LLC tables never
invalidate what an older parked read still needs. In the \emph{pipelined} model
(\key{m\_data\_handler.m\_data\_array\_pipelined}\code{=1}) every access pays the latency and one
is admitted per port per cycle (\key{m\_data\_handler.m\_data\_array\_ports}); a message whose
row touches the array is deferred \emph{whole}, so that its state change and its data access
happen together when the access runs. Lines that sit in the MSHR or the write-back buffer are
not in the array and bypass the pipeline in either model.

Deferring a message whole takes it out of the queue, so a younger message to the same line could
otherwise act on a state the older one was about to change. The \emph{line interlock}
(\key{line\_interlock}\code{=1}) prevents that: while a line has an array access pending, every
queued message to it is held in place, not ready, until the access has completed, so each message
sees the line in bus order. The pipelined model turns the interlock on; with real data flowing
through the hierarchy (the gem5 integration, \cref{ch:integration-modes}) or an L1 with an array
latency, it must be on. In the same spirit a message from the bus enters a controller's queue
ahead of the core's own requests but behind older bus messages, so bus order is kept even when a
bus message has to wait. \file{docs/StateAndData.md} walks through the situations in which line
state and line data can come apart and which setting handles each.

\begin{table}[htbp]\centering\small
\caption{The data-array settings.}\label{tab:array-models}
\begin{tabularx}{\textwidth}{@{}lllY@{}}
\toprule setting & default & gem5 preset & effect \\ \midrule
\code{m\_data\_array\_pipelined} & 0 & 1 (LLC) & accesses pay the latency, one admitted per port per cycle; always defers whole, so it turns the interlock on \\
\code{m\_data\_array\_ports} & 1 & 1 & pipelined model only: accesses admitted per cycle \\
\code{line\_interlock} & 0 & 1 (L1, LLC) & hold every message to a line with an array access pending; enables whole-message deferral in the occupancy model \\
\code{m\_data\_access\_latency} & 0 (L1), 10 (LLC) & 0, 10 & the port's busy time per access \\
\bottomrule
\end{tabularx}
\end{table}

\section{Admission and back-pressure}

A new demand miss is admitted only when the controller can finish it. \code{canAdmitRequest}
stalls a demand request at the lower interface when the MSHRs are full or the write-back buffer
has no headroom; \code{demandAdmissionBlocked} adds those two bounds to the processing-queue
bound. Responses and the controller's own maintenance traffic are never subject to these
checks, which is what keeps the system deadlock-free under finite buffers: a full queue of
stalled requests can always drain, because the data they wait for is always admitted.

The default depths of every queue and buffer are in \cref{app:config-keys}.

\section{The LLC's extra duties}

The shipped LLC (\code{CacheController\_End2End}) is \emph{inclusive}: every line an L1 holds is
also in the LLC. That gives it two duties an L1 does not have. When it evicts a line it must
remove the L1 copies, which it does with the FSM action \act{IssueInv}: an \code{INV} broadcast on
the TripleBus \emph{service} lane that every L1 sees and that comes back to the LLC itself as
\code{Own\_Invalidation} to complete the eviction (\cref{ch:llc-datapath}). And it tracks, per
line, which L1 currently owns it (\act{SetOwner}/\act{ClearOwner}), so that it can stay out of the
way when an owner answers a request on the bus. The directory controller
(\code{CacheControllerDirectory}) replaces snooping with a sharer list at the home and forwards
requests point-to-point; \cref{ch:coherence-engine} contrasts the two.

\section{Message encoding}\label{sec:message-encoding}

Every transaction is one \code{Message} object. The decoders in the protocol handlers read a
small set of fields (\cref{tab:message-fields}); knowing them is what makes a Debugger trace or
the raw event trace readable.

\begin{table}[htbp]\centering\small
\caption{The \code{Message} fields a decoder reads.}\label{tab:message-fields}
\begin{tabularx}{\textwidth}{@{}>{\ttfamily}lY@{}}
\toprule \normalfont field & meaning \\ \midrule
source & where the message entered the controller: \code{LOWER\_INTERCONNECT} (from the cores' side), \code{UPPER\_INTERCONNECT} (from the bus or memory side), \code{SELF} (a replacement the controller raised for itself) \\
data & \code{NULL} for a request or command; non-\code{NULL} for a data transfer \\
complementary\_value & on a request, its type: 0 GetS, 1 GetM, 2 PutM, 10 INV (CPU$\rightarrow$L1: 0 Load, 1 Store). On data it is \emph{inherited} from the message the data was copied from: provenance, not a kind \\
owner & the requester on requests, supplies and fills; the evicting L1 on an eviction write-back; the LLC on an invalidation \\
to & destination interface ids; a response is delivered only to these, a broadcast request to everyone \\
msg\_id & a demand request keeps one id from \code{GetS}/\code{GetM} through fill and response; an eviction write-back has its own id; a write-back forced by a back-invalidation carries the LLC eviction's id \\
kind & the producer's classification of the message (\code{DEMAND}, \code{GETS}, \code{GETM}, \code{PUTM}, \code{INV}, \code{MEM\_READ}, \code{EVICT}, \code{WB\_DATA}, \code{WB\_INV}, \code{SUPPLY}, \code{SUPPLY\_DEFERRED}, \code{RESP}, \code{FILL}, \code{FILL\_ROLLBACK}, \code{MEM\_WRITE}), set where it is created because the fields above cannot be read back into one \\
\bottomrule
\end{tabularx}
\end{table}

The bus \code{MessageType} passed to \code{pushMessage} (\code{REQUEST}, \code{DATA\_RESPONSE},
\code{SERVICE\_REQUEST}) selects the lane the message travels on; it is not stored in the message.
Three things the same value means are worth remembering: \code{cv = 2} is a PutM on a request, an
eviction write-back at the LLC and an exclusive grant at a MESI L1; \code{owner} is the requester
for requests, supplies and fills but the LLC for invalidations; and \code{msg\_id} joins a demand
request across its life, while an invalidation-forced write-back is attributed to the LLC
eviction, not to any core's request. The full producer-by-producer table is in
\file{docs/MessageEncoding.md}.

\section{Delivery order at the bus interface}

A snoop protocol is correct only if every controller observes bus transactions in the
\emph{same} order. \code{BusInterface} therefore stamps each delivery into a controller's RX
with one global sequence number, and \code{TripleBusInterface} hands out a service-lane message
(an LLC back-invalidation) only when it was delivered \emph{before} the request at the head of
the request RX. Data responses are deliberately not ordered against requests, so that a full
queue of stalled requests can always drain.

\begin{wherebox}
\begin{itemize}[nosep,leftmargin=1.2em]
\item \file{header/CacheControllers/BaseController.h}, \file{src/CacheControllers/BaseController.cpp} --- the machinery every controller shares
\item \file{header/CacheControllers/CacheController.h}, \file{src/CacheControllers/CacheController.cpp} --- the data array, miss handling, write-back path
\item \file{header/CacheControllers/CacheControllerExclusive.h}, \file{src/CacheControllers/CacheControllerExclusive.cpp} --- the shipped L1
\item \file{header/CacheControllers/CacheController_End2End.h}, \file{src/CacheControllers/CacheController_End2End.cpp} --- the shipped LLC
\item \file{header/CacheControllers/CacheControllerDirectory.h}, \file{src/CacheControllers/CacheControllerDirectory.cpp} --- the directory home
\item \file{src/CacheDataHandler_COTS.cpp} --- array, MSHRs, write-back buffer, port arbiter
\item \file{header/FRFCFS_Buffer.h} --- the processing queue
\item \file{src/Protocols/} --- the handlers (\code{readEvent}, \code{handleAction})
\item \file{header/Interconnect/CommunicationInterface.h} --- the message
\item \file{src/Interconnect/BusInterface.cpp}, \file{src/Interconnect/TripleBusInterface.cpp} --- delivery order
\end{itemize}
\end{wherebox}
